%!TEX root = ../atomic_jsc.tex
% LTeX: language=en
\section{Hilbert's basis property}
%
A monoid action $\monoid\actson\Indets$ is said to satisfy the \intro{Hilbert's basis property} if every \kl{$\monoid$-ideal} in $\poly{\K}{\Indets}$ is \kl{orbit-finitely generated},
equivalently,
every strictly increasing sequence of $\monoid$-ideals is finite.
The equivalence of two definitions can be proven by a minor modification of the proof of equivalence of two definitions of Noetherianity \arka{cite wiki}.
The reader can find a proof in the \arka{move the proof to appendix and refer to it}.
%
\begin{proof}[Proof of equivalence of two definitions of the Hilbert's basis property]
%
(orbit-finite generation $\implies$ finiteness of strictly increasing chains)
Take an infinite strictly increasing sequence $\idl_1\subsetneq\idl_2\subsetneq \dots$ of $\monoid$-ideals.
Let $\idl[J]$ be the ideal generated by $\bigcup_{i=1}^n \idl_i$.
Because each $\idl_i$ is $\monoid$-equivariant, $\idl[J]$ is an $\monoid$-ideal.
Let $F$ be an orbit-finite set of generators of $\idl[J]$.
By definition of orbit-finitness There exists $f_1,\dots,f_k\in F$ such that $F = \bigcup_{j=1}^k \orbit[\monoid]{f_j}$.
For every $j\in\{1,\dots,k\}$, there exists $i_j\in\N$ such that $f_j\in\idl_{i_j}$.
Since the sequence of $\idl_j$ is increasing,
$f_1,\dots,f_k\in\idl_{i_k}$.
Due to $\monoid$-equivariance of $\idl_{i_k}$,
$F\subseteq \idl_{i_k}$.
But then $\idl_{i_k + 1} \subseteq \idl[J]\subseteq \idl_{i_k}$ which contradicts the fact that $\idl_{i_k}\subsetneq \idl_{i_k + 1}$.

(finiteness of strictly increasing chains $\implies$ orbit-finite generation)
Let $\idl[j]$ be an $\monoid$-ideal.
Iteratively pick $f_0,f_1,f_2,\dots\in\idl[J]$ such that $f_0\in\idl[J]$ and $f_{k+1}$ is not in the ideal generated by $\idl_k = \bigcup_{i=1}^k \orbit[\monoid]{f_i}$ for every $k\in\N$.
We have a strictly increasing chain $\idl_1 \subsetneq \idl_2 \dots$,
which must be finite by assumption.
Let $\idl_n$ be the last element in this chain.
We have $\idl[J]\setminus\idl_n = \emptyset$.
By construction $\idl_n\subseteq \idl[J]$.
Hence $\idl_n = \idl[J]$.
But then $\bigcup_{i=1}^k \orbit[\monoid]{f_i}$ is an orbit-finite set of generators of $\idl[J]$.
%
\end{proof}
%
The papers \cite{LauSno26,GHOLAS24} give necessary and sufficient conditions for an action $\monoid\actson\Indets$ to have the Hilbert's basis property.
Below we describe these conditions.

Define the \intro{divisibility up-to $\monoid$} relation ($\intro*\gdivleq[\monoid]$) as follows:
for $\monelt[p], \monelt[q] \in\mon{\X}$, we write $\monelt[p] \gdivleq \monelt[q]$ if there exists $\gelem \in \monoid$ such that $\gelem \cdot \monelt[p]$ \kl{divides} $\monelt[q]$.
We say that a monoid action $\monoid \actson \Indets$ is \kl{nice} if $\gdivleq[\monoid]$ is a \kl{well-quasi-order} on $\mon{\Indets}$.
%
\begin{theorem}[{\cite[Theorem 11]{GHOLAS24}}]
If a monoid action has the Hilbert's basis property then it is \kl{nice}.
\end{theorem}
%
The above theorem can be proven by generalising some of the ideas used in the proof of the fact that the ideal $\calC$ defined in \Cref{ex:cycles} is not orbit-finitely generated.
The ideas used to prove this theorem will be useful in many of our proofs.
Thus, we do a quick recall of it using our notation.
%
\begin{proof}
Consider an infinite sequence $\monelt_0,\monelt_1,\monelt_2,\dots$ of monomials in $\mon{\Indets}$.
We show there exists $m < n\in\N$ such that $\monelt_m\gdivleq[\monoid]\monelt_n$.
Let $\idl[M]_k$ be the ideal generated by $\bigcup_{i=0}^k \orbit[\monoid]{\monelt_i}$.
Clearly, $\idl[M]_k \subseteq \idl_{k+1}$ for all $k\in\N$.
Pick $n \geq 1$ such that $\idl[M]_{n-1} = \idl[M]_{n}$.
Such an $n$ must exists because otherwise we have a strictly increasing chain of $\monoid$-ideals $\idl[M]_0\subsetneq\idl[M]_1\subsetneq\idl[M]_2\dots$ contradicting the Hilbert's basis property of $\monoid\actson\Indets$.
Since $\idl[M]_{n-1} = \idl[M]_{n}$, $\monelt_{n}\in\idl_{n-1}$.
\arka{put a remark that $0$ is not a monomial}
This implies there exists $r_1,\dots,r_k\in\K\setminus\{0\}$, $\monelt[p]_1,\dots,\monelt[p]_k\in\mon{\Indets}$ and $\monelt[q]_1,\dots,\monelt[q]_k\in\bigcup_{i=0}^{n-1} \orbit[\monoid]{\monelt_i}$ such that
\[
\monelt_{n} = \sum_{i = 1}^{n-1} r_i\monelt[p]_i\monelt[q]_i \ .
\]
The equality can hold only if $\monelt[q]_{i'}$ divides $\monelt_{n}$ for some $i'\in\{0,\dots,n-1\}$.
Let $m \in\{0,\dots,n-1\}$ be such that $\monelt[q]_{i'}\in\orbit_[\monoid]{\monelt_m}$.
But then $m < n$ and $\monelt_m\gdivleq[\monoid]\monelt_n$.
\end{proof}
%
A monoid action $\monoid \actson \Indets$ is called \intro{ordered} if it respects a total order $\varlt$ on $\Indets$ (i.e., $\gelem \cdot x \prec \gelem \cdot y$ for every $x \varlt y$ and $\gelem\in\monoid$).
It is \intro{nicely ordered} if it is \kl{nice} and \kl{ordered}.
%
\begin{theorem}[{\cite[Theorem 12]{GHOLAS24}, \cite[Proposition 3.2]{LauSno26}}]
\label{thm:ofgen}
If a monoid action is \kl{nicely ordered} then it has the \kl{Hilbert's basis property}.
\end{theorem}
%
We skip the proof of this theorem as we prove a stronger statement later \arka{refer to theorem ???}.
Theorem ??? says that if a monoid action $\monoid\actson\Indets$ is \kl{nicely ordered} then every $\monoid$-ideal has an orbit-finite $\monoid$-Gr\"{o}bner basis.
Since a $\monoid$-Gr\"{o}bner basis of an ideal generates the ideal \arka{cite lemma},
this implies every $\monoid$-ideal has an orbit-finite set of generators whenever $\monoid\actson\Indets$ is \kl{nicely ordered}.
%
A \kl{monoid action} $\monoid\actson\Indets$ is called a \intro{reduct} of another \kl{monoid action} $\mathcal{N}\actson\Indets$ if for every finite subset $S$ of $\Indets$ and $\gelem\in\mathcal{N}$,
there exists $\sigma\in\monoid$ such that $\pi\cdot x = x$ for every $x\in S$.
This implies that every $\mathcal{N}$-equivariant subset is also $\monoid$-equivariant.

Reducts of \kl{nicely ordered} \kl{monoid actions} are called \intro{nicely orderable}.
Observe that a \kl{nicely orderedable} monoid action must be \kl{nice},
however it may not be \kl{nicely ordered}.
For instance, $\symgr{\Indets}\actson\Indets$ is \kl{nicely orderable} but is clearly not \kl{nicely ordered} as it is not \kl{ordered}.

Applying this observation and using \Cref{thm:ofgen} we get:
%
\begin{corollary}[{\cite[Corollary 13]{GHOLAS24}}]
If a monoid action is \kl{nicely orderable} then it has the \kl{Hilbert's basis property}.
\end{corollary}
%
At this point we must state Pouzet's conjecture.
%
\begin{conjecture}[{\cite[Problems 12]{POUZ24}}]
Every \kl{nice} monoid action is \kl{nicely orderable}.
\end{conjecture}
%
\begin{remark}
\arka{Conjecture fails without injectivity}
\end{remark}
%
Our reformulation of this conjecture is slightly different from the original one.
In \arka{refer to sec} we will clarify how these two formulation are (almost) equivalent. 

Observe that a positive answer to this conjecture will prove that a \kl{monoid action} has the Hilbert's basis property if and only if it is \kl{nice}.